\section{Discussion}
\label{sec:discussion}

\subsection{Attacking the Metaphor, Deliberately}
\label{sec:disc-metaphor}

A concept paper owes its readers the strongest case against itself. We therefore
attack our own framing and report what survives. \emph{Where the financial analogy fails:} there is no
creditor and no contract---nobody agreed to the terms, and the ``lender'' (future
circumstance) never negotiates; effort and abstention exposure under sound exhaustive
policies follow a stochastic
schedule of cliffs and hazards while error may drift with discriminator quality
(\cref{prop:steps}), so compounding
intuitions---and any metric built on them---mislead (\cref{met:interest} was
rejected on exactly this ground); policy burden is not fungible---effort, accepted
error, and abstention (\cref{fnd:effort,fnd:false,fnd:redundancy}) are different
channels that do
not net against each other; and physical irrecoverability boundaries are mandatory and world-imposed
(\cref{prop:writeoff}), where financial write-offs are creditor choices.
\emph{Where the analogy earns its keep:} the deferral structure is real (a present
choice with a later, larger, contingent cost); the principal/carrying-cost
decomposition marks a real operational distinction (capture-time gaps versus
decay-driven growth) with distinct remedies (capture discipline versus retention
and export policy); and ``debt'' correctly imports the management insight that the
stock must be measured and budgeted, not merely deplored.

\emph{The removal test}: delete the word ``debt'' and every construct still
stands---deficit (\cref{def:deficit}) is a syntactic census; reconstruction cost
(\cref{def:rec}) is a measured quantity; decay schedules and half-lives
(\cref{sec:interest}) are computed from published retention and turnover;
irrecoverability (\cref{def:irrec}) is a decidable/certifiable state; and the
expected-excess-burden aggregate (\cref{def:debt}) is an ordinary expected value that
could as well be called \emph{reconstruction exposure}. We embrace that reading:
the \emph{measured} quantity is reconstruction exposure, and a skeptic may use
that vocabulary throughout without loss of formal content.

What, then, does the debt framing add that exposure does not? Exposure language
frames the quantity as environmental---a risk one has, like weather. The debt
framing makes one further, contentful commitment: it binds the aggregate to
\emph{identifiable, dated, attributable deferral decisions}---this change, shipped
ticketless, under this pressure, by this team---and thereby licenses management
machinery that exposure language does not naturally carry: per-decision
registration at creation time (the register of \cref{sec:disc-management}, with
backfill due dates set inside the cheapest-path window); attribution of the
aggregate to controllable practices P1--P8 rather than to fate; and the
inflow/stock separation (\cref{met:growth} versus interest) that tells an
organization whether it is borrowing faster or merely aging. These are decisions
and accountabilities, not measurements; the measurement structure is
metaphor-free, and the metaphor's residual work is to keep the human origin of
the exposure attached to the number. That is a naming choice we defend as
useful---while conceding, per the analysis above, that every quantitative
intuition imported from finance (rates, compounding, refinancing) must be checked
against the model and mostly fails.

Circularity was the other standing risk, and the definitional order of
\cref{sec:definitions} was built against it: deficit is defined without reference
to cost (syntax against a declared schema); cost is defined without reference to
deficit (measured effort of answering queries); debt relates the two through an
explicit counterfactual; and the schema- and workload-relativity are declared
inputs rather than smuggled conclusions. What \emph{is} irreducibly judgmental
---the choice of $\Sch$ and $(w_q, \lambda_q, \pi_q)$---is surfaced as a governance decision,
not hidden in a formula (\cref{sec:model-notformula}).

\subsection{Relativity of the Construct}
\label{sec:disc-relativity}
The sharpest structural limitation is not the metaphor but the relativity:
$\ED$ is defined against a declared schema $\Sch$ and a declared workload
$(w_q, \lambda_q, \pi_q)$, all controlled by the party being measured. An estate can reach
$\ED \approx 0$ by declaring a weak schema or a sparse workload, and
\cref{fnd:ed} showed concretely how the penalty declaration reshapes what the
number says. We therefore state plainly: what is measurable
\emph{inter-subjectively} is the ingredient vector---EDR against a published
schema, measured ERC and FRR under the same stated policy, the decay schedule, and the
irrecoverability constructs---while $\ED$ itself is an accounting quantity
relative to published governance inputs, comparable across estates only under
shared reference declarations and the same $\rho$. Policy choice is another gaming
surface: an aggressive policy can trade $N$ for $W$, while a conservative one does
the reverse. Cross-organizational reporting should therefore publish S/W/N and IRR,
and, where policies differ, report a policy frontier rather than compare scalar
$\ED_\rho$ values as though they shared an estimand. To make comparison possible rather than merely
honest, we propose (as future standardization, not as this paper's result) a
\emph{reference declaration}: the six-interrogative schema of
\cref{sec:schema-informal} as $\Sch_{\mathrm{ref}}$, and a workload
$\Qw_{\mathrm{ref}}$ parameterized by public covariates---per-service audit
queries at a fixed annual rate, incident queries at the estate's measured
incident rate, one recovery query per service---with error and abstention weights published
alongside any reported figure. This is the same discipline financial accounting
imposes through reporting standards: the number is meaningful because the rules
for producing it are shared, not because it is observer-independent.
For initial calibration, teams should derive $w_q$ from observed incident/audit
query counts per reporting horizon or from explicit scenario probabilities;
$\lambda_q$ from the incremental consequence of acting on an accepted wrong answer
(e.g., investigation labor, remediation, or outage minutes); and $\pi_q$ from the
incremental consequence of leaving that same question unresolved. All three must be
expressed in one declared accounting unit. A practical elicitation is to record
low/base/high values with the incident commander, audit owner, and service owner,
then publish both the elicited rationale and the resulting grid. No universal
DORA-derived default exists for these losses; when local estimates are weak, the
honest output is the ingredient vector plus a declared $\lambda/\pi$ grid such as
\cref{tab:pi-sensitivity}, not a falsely precise scalar. IRR remains a separate
physical diagnostic unless a joint loss $L_q(Y_\rho,\mathsf I_q)$ is calibrated.

\subsection{Explanatory and Predictive Value}
The concept must be useful beyond naming. Its explanatory content: it unifies
under one mechanism phenomena currently discussed separately---why incidents on
old systems take longer to diagnose (interest), why audits of healthy systems fail
(deficit despite health), why tool migrations feel like organizational amnesia
(bulk path death), why the departure of one engineer disproportionately damages
some services (human path as last survivor). Its predictive content is concrete
and falsifiable: effort plateaus and cliffs aligned to retention/migration schedules (IH1),
departure-boundary cost multiples (IH2), memory as modal final path (IH3), and
the simulation-grounded expectation that effort inflation precedes failure
(\cref{fnd:effort}). Accepted wrong answers occur, but the isolated density trend
is mixed (\cref{fnd:false}). \Cref{sec:field} specifies the refutation
conditions; a concept that survives them will have earned its place, and one that
does not will have been usefully eliminated.

\subsection{Management Implications (Proposals)}
\label{sec:disc-management}
Stated as proposals, not findings. \emph{Evidence-debt register}: like a TD
register~\cite{seaman2011,ernst2015}, an explicit ledger of known deficits with
principal estimates and decay deadlines---P2 break-glass changes, for instance,
entered at creation with a backfill due date inside the cheapest-path window
(\cref{sec:interest-worked}). \emph{Retention as governance, not default}:
\cref{sec:interest-halflife}'s observation that half-lives are set by unexamined
vendor defaults makes retention review the highest-leverage, lowest-cost
intervention. \emph{Backfill sprints as repayment}: deliberate reconstruction of
high-value history while paths are cheap---though \cref{fnd:false} warns that late
backfill in dense estates risks writing confident fiction into the record, so
backfilled facts should carry confidence annotations. \emph{Deferral budgets}:
if evidence capture is to be skipped under pressure, the skip should be a recorded
decision---the debt family's central lesson~\cite{kruchten2012} transposed: unmanaged
debt is not the debt you took on but the debt you never knew you took on.
\emph{Capture-versus-risk gate}: for capture action $a$ over horizon $H$, record
\[
K_H(a)=K_{\mathrm{labor}}(a)+K_{\mathrm{storage}}(a)
       +K_{\mathrm{privacy/control}}(a),
\]
and adopt $a$ when
\[
K_H(a)<\sum_s p_s\bigl[\ED_{\rho,s}^{\neg a}
                         -\ED_{\rho,s}^{a}\bigr],
\]
with scenario probabilities, horizon, accounting unit, and uncertainty range
published. Otherwise selective omission may be rational.
The comparison must be recorded because high-volume controllers can make exhaustive
capture technically or legally disproportionate. This is a decision rule proposal,
not a claim that this paper has measured either side of the trade-off; it closes the
decision logic without pretending that the synthetic study supplies capture costs.

\subsection{Limitations}
Beyond the simulation threats (\cref{sec:sim-threats}): the framework prices
reconstruction, not the first-order value of evidence for real-time operations
(alerting, debugging), which observability already treats~\cite{majors2022}; the
workload $\Qw$ is a declared judgment and organizations may declare it
self-servingly---the discipline is publishing it; capture itself has costs
(effort, storage, privacy exposure under data-protection law~\cite{gdpr,
nist80092}) that our gate parameterizes but this study does not calibrate, so the
empirical optimization of capture under cost and legal constraint remains open;
and the only non-synthetic observation is a bounded, age-stratified public-forge
sample that shows missing external issue linkage occurs under a declared schema
but measures neither intent adequacy, decay, reconstruction, nor money
(\cref{tab:argocd-census}). All cost, error, interaction, and debt magnitudes remain
synthetic, a boundary marked at every claim.
